%% kapa: exam=1 round=33 year=2022 period=2 form=A subject=appraisal_law subject_ko=감정평가관계법규 seq=20 number=20
%% source: Q-Net 감정평가사 기출문제 (cst00309) artlSeq=5209997 file=raw/1008/5209997/2022년도 제33회 감정평가사 1차시험 2교시(감정평가 관계법규, 회계학).pdf
%% conversion: kordoc 4.24.51 -> tools/gen1.py (auto); review=proofread-sonnet; flags=none
\begin{question}{20}
국유재산법령상 국유재산에 관한 설명으로 옳은 것은?
\begin{choices}
  \item 국가가 직접 사무용·사업용으로 사용하는 재산은 공공용재산이다.
  \item 총괄청은 일반재산을 보존용재산으로 전환하여 관리할 수 있다.
  \item 중앙관서의 장등이 필요하다고 인정하는 경우에는 보존용재산에 사권을 설정할 수 있다.
  \item 공용재산은 시효취득의 대상이 될 수 있다.
  \item 영농을 목적으로 하는 토지와 그 정착물의 대부기간은 20년 이내로 한다.
\end{choices}
\end{question}
\answer{2}
